\chapter*{Abstract}
\addcontentsline{toc}{chapter}{Abstract}

Large Language Models (LLMs) are increasingly used in coding assistants and agentic workflows, where small prompt-framing choices affect parser compatibility, verbosity, latency, and correctness. However, common advice about tone, role assignment, and formats like JSON or XML remains under-tested across models and tasks.

This thesis examines prompt-framing effects through a full-factorial experiment with 18 prompt variants (3 tones $\times$ 2 roles $\times$ 3 formats). With 50 base tasks per task kind (Python code generation and mathematical reasoning), we conducted 1,800 runs per model across 6 local and 2 cloud LLMs, resulting in 14,400 runs in total. We measure Pass@1 correctness, duration, completion tokens, and output compliance.

Results show that model ability dominates correctness, while prompt effects are smaller, task-specific, and model-specific. Nemotron and DeepSeek strongly skew pooled averages. A mixed-effects logistic analysis with base-task random intercepts confirms the main contrasts. Prompt framing affects behaviors more than correctness: polite prompts increase output length, expert role framing slightly hurts cloud model performance, and JSON reduces the math pass rate for local models. In contrast, structured formats (JSON/XML) improve pass rates in API Python code generation.

The study found no universal prompting technique that reliably improves model performance. A strategy that helps one model may hurt another. Therefore, prompt engineers should test policies on their target models, instead of using universal prompt advice. Reliable LLM-based coding systems need explicit output constraints, robust parsers, validation loops, and local prompt-policy testing.

\vspace{1cm}
\noindent\textbf{Keywords:} Large Language Models, Prompt Engineering, Code Generation, Mathematical Reasoning, Empirical Software Engineering
